\documentclass[11pt]{article}
\usepackage[margin=2.5cm]{geometry}
\usepackage{amsmath,amssymb}
\usepackage{booktabs}
\usepackage{xcolor}
\usepackage{hyperref}

% placeholders: current numbers, to be replaced by tables generated from the workflow
\newcommand{\cur}[1]{\textcolor{blue}{[#1]}}
\newcommand{\todo}[1]{\textcolor{red}{[TODO: #1]}}

\title{A learned prior over model forms for biology and ecology\\[0.3em]
\large Paper 1: outline (draft)}
\author{}
\date{}

\begin{document}
\maketitle

\noindent Blue brackets are current numbers from the workflow (\texttt{results/grammar/*.md}); red brackets
are to-dos. Every table should be generated by a Snakemake target and \verb|\input| here, not typed.

%=========================================================================
\section*{Working abstract}
Choosing the functional form of a model is usually done by convention. Symbolic regression and Bayesian
model selection both need a prior over forms, and the prior matters most exactly when data are weak.
We assemble a corpus of \cur{$\sim$40k model--form pairs from 12 sources} of published biological and
ecological models, define a family of proper priors over expression trees (from symbol frequencies to
hierarchical Pitman--Yor context models and adaptor grammars that reuse whole forms and subtrees), and
test them against four kinds of forms they were not trained on: whole held-out sources, forms first
published after a cutoff year, curated catalogues never used for training, and forms derived from random
mechanisms rather than from the literature. Learned priors give unseen forms \cur{7--8 nats} more log
probability than a symbol-frequency baseline, and the adaptor priors also predict forms that appear
later in the literature best. A look-alike test shows \todo{headline}.

%=========================================================================
\section{Introduction}
\begin{itemize}
  \item Form choice by convention (Michaelis--Menten, logistic, Arrhenius, \dots); how this interacts
  with weak data and model selection.
  \item Symbolic regression needs a prior over expressions; existing ones: implicit (search
  heuristics, complexity penalties), uniform over a grammar, or learned from physics equations
  (Bartlett, Desmond \& Ferreira 2023, \texttt{fnprior}; Feynman equations). None for biology/ecology.
  \item Questions: (1) Can a prior learned from the literature generalise to forms it has not seen?
  (2) Is it just recapitulating convention? (3) Does it help select the true form when data are weak?
  \item Contributions: corpus and pipeline; family of proper tree priors with a clean baseline; an
  evaluation protocol (raw proper log P; seen/novel; size-matched; four kinds of out-of-sample test;
  look-alike test); recommendations.
\end{itemize}

%=========================================================================
\section{A corpus of model forms}
\begin{itemize}
  \item Sources (training, 12): BioModels curated and non-curated, CellML, ModelDB (NMODL), SABIO-RK,
  JWS Online, curvelib, nlmixr2lib, allodb, thermal performance, ecological curves, BioCat. Table:
  models, model--form pairs, distinct forms, licence, extraction route.
  \item Hold-out sources (never trained on): rnls (formulas fitted in CRAN package code),
  ratelaws (SBO, COPASI), ecocat2 (mortality, abundance, P--I, soil respiration, stomatal,
  transmission, dose--response, phenology). \cur{110 / 344 / 89 pairs; 14 / 1 / 32 novel forms}.
  \item From expressions to forms: canonicalisation (simplify, order, rename constants);
  generalisation of numerical literals; the redundant-constant check (Jacobian rank).
  \todo{report after the canonicaliser fix: deterministic rebuild, merged redundant forms}.
  \item Overlap between sources (Venn/upset figure); size distribution; most frequent forms per source.
  \item Dates (for the temporal split): coverage per source \cur{BioModels, ModelDB, SABIO-RK complete;
  JWS 78\%; CellML heuristic}.
  \item Reproducibility: Snakemake workflow, pinned environments, cached raw downloads with manifests.
\end{itemize}

%=========================================================================
\section{Priors over model forms}
Main text: a condensed version of \texttt{docs/prior\_models.tex}; full version as supplement.
\begin{itemize}
  \item Forms as trees; product of per-node choices given a context; properness (branching-process
  argument; uniform symbols is improper).
  \item Hierarchical Pitman--Yor over contexts; inference; Kneser--Ney (\texttt{fnprior}) as its
  fixed-parameter limit.
  \item Variants: \texttt{sym0} baseline, \texttt{hpy2}, \texttt{hpy3\_depth}, \texttt{hpy2\_src},
  adaptors (\texttt{form\_ad}, \texttt{form\_ad3\_depth}, \texttt{tree\_ad3}); PCFG grammars as an
  alternative family.
  \item Blended prior $w P_{\text{adaptor}} + (1-w) P_{\text{context}}$, $w$ chosen by leave-one-source-out.
  \todo{results of the blend target}.
  \item Training weightings: pooled vs equal per source.
\end{itemize}

%=========================================================================
\section{Evaluation protocol}
\begin{itemize}
  \item Score: raw log P per test model--form pair (proper; comparable across tree priors). Why not
  renormalise over a candidate list (rewards flatness over the list).
  \item Seen vs novel test forms; size-matched comparisons (OLS novelty penalty at fixed node count;
  gain over \texttt{sym0}); redundant forms excluded as a check.
  \item Four out-of-sample designs, each answering a different criticism:
  \begin{enumerate}
    \item leave one source out (does it transfer across communities?);
    \item hold-out sources (curated catalogues never trained on);
    \item temporal split (does it predict forms the field adopts later?);
    \item mechanism-derived forms (is it more than convention?).
  \end{enumerate}
  \item The look-alike test: simulate data from a true form; candidates whose evidence is within
  $\Delta$ nats of the best; does the prior rank the truth first among them?
\end{itemize}

%=========================================================================
\section{Results}

\subsection{Learned priors beat the baseline, on seen and novel forms}
\begin{itemize}
  \item LOSO table (per source; equal/pooled). Adaptors best overall \cur{1--4 nats over
  \texttt{fnprior}}; context models best on novel forms in LOSO.
  \item Hold-out sources: \cur{adaptor $-6.3$, \texttt{hpy3\_depth} $-7.0$, \texttt{fnprior} $-7.9$,
  \texttt{sym0} $-15.2$ nats per pair}. PCFG grammars do not beat \texttt{sym0}.
\end{itemize}

\subsection{Why novel forms score low: size, symbols, structure}
\begin{itemize}
  \item Seen--novel gap \cur{$\approx$17 nats for \texttt{fnprior}}: \cur{$\approx$10} from size,
  \cur{$\approx$2} from rarer symbols, \cur{4--5} from unfamiliar structure; the structural part
  is larger for priors that rely more on seen forms (\cur{$\approx$13 for adaptors}).
  \item Gain over \texttt{sym0} is as large on novel forms as on seen ones for the context models
  (\cur{$+7.2$ / $+7.6$}): they generalise rather than memorise.
  \item Figure: log P vs node count, seen and novel, per prior.
  \item \todo{same diagnostics for the temporal split}.
\end{itemize}

\subsection{Predicting the future literature}
\begin{itemize}
  \item Temporal split, cutoffs 2010/2015/2020: \texttt{tree\_ad3} best on all, seen and novel
  forms in all six settings (\cur{novel $-28$ to $-31$ vs $-36$ to $-39$ for \texttt{fnprior}}).
  \item New forms are recombinations of existing subtrees; contrast with hold-out catalogues, where
  novel forms are structurally new and adaptors lose.
  \item Caveat: ModelDB dominates the test sets. \todo{ModelDB-excluded table}.
\end{itemize}

\subsection{Is the prior just convention? Mechanism-derived forms}
\begin{itemize}
  \item Random King--Altman enzyme mechanisms: \cur{12 distinct forms from 1959 mechanisms; the 8 of
  degree $\le 2$ are all in the corpus; 4 cubic-denominator forms are not}. Conventions match
  mechanisms at low complexity; frequencies differ (substrate inhibition common mechanistically,
  rare in the corpus).
  \item \todo{compartment time courses and temperature dependence; prior scores on mechanism forms,
  seen/novel}.
\end{itemize}

\subsection{Does the prior pick the right form when data are weak? The look-alike test}
\begin{itemize}
  \item \todo{table: per prior, rank of truth within look-alike set, renormalised prior mass on
  truth, MAP accuracy; by seen/novel, mechanism class, $(n, \sigma)$}.
  \item The price of a prior that bets on the literature: how often a novel truth is overridden.
\end{itemize}

\subsection{Sensitivity}
\begin{itemize}
  \item Equal vs pooled weighting; source-labelled hierarchy (no gain).
  \item \todo{prior tempering: $P^{\tau}$ and the look-alike results as $\tau$ varies}.
  \item \todo{effect of the canonicaliser fix on all headline numbers}.
\end{itemize}

%=========================================================================
\section{Discussion}
\begin{itemize}
  \item Recommendation: which prior to use, and when (blend; or adaptor if forms are expected to be
  conventional, context model if not).
  \item What a prior over forms can and cannot do; complementary to data; reporting the prior's
  contribution to a selected form (the evidence decomposition, used in paper 2).
  \item Limitations: one variable; forms only (not constants); canonicalisation choices; corpus
  coverage (neuroscience-heavy ModelDB, physiology-heavy CellML); dated sources only partly.
  \item Link to paper 2 (hierarchical symbolic regression with this prior as the base measure).
\end{itemize}

%=========================================================================
\section*{Supplement}
\begin{itemize}
  \item Full description of every prior (\texttt{docs/prior\_models.tex}).
  \item Source extraction details and licences (\texttt{hsr/data/sources/README.md}).
  \item Properness checks; MCMC diagnostics for the HPY and adaptor fits.
  \item All per-source tables.
\end{itemize}

\end{document}
